\documentclass[a4paper,11pt]{article}

\usepackage{fullpage}
\usepackage[all]{xy}
\usepackage[utf8]{inputenc}
\usepackage{amsmath,amsthm}
\usepackage{amsfonts}
\usepackage{graphicx}
\usepackage{latexsym}
\usepackage{float}
\usepackage{color}
%%%%%%%%%%%%%%% COMMANDS %%%%%%%%%%%%%%%%%%%%%%%%%%%

\newcommand{\itemb}{\item[$\bullet$]}

\newcommand{\Fd}{\mathbb{F}}
\newcommand{\Znat}{\mathbb{Z}}
\newcommand{\Nnat}{\mathbb{N}}

%%%%%%%%%%%%%%% ENVIRONMENTS %%%%%%%%%%%%%%%%%%%%%%%

\theoremstyle{plain}
\newtheorem{lemme}{Lemma}
\newtheorem{algorithme}{Algorithm}
\newtheorem{corolaire}{Corollary}
\newtheorem{propriete}{Property}
\newtheorem{proposition}{Proposition}
\newtheorem{theoreme}{Theorem}
\newtheorem{definition}{Definition}

\theoremstyle{definition}
\newtheorem{exercice}{Exercise}
\newtheorem{remarque}{Remark}
\newtheorem{exemple}{Example}


%%%%%%%%%%%%%% DOCUMENT %%%%%%%%%%%%%%%%%%%%%%%%%%%%


\begin{document}

\hbox{\raisebox{0.4em}{\vrule depth 0pt height 0.5pt width \textwidth}}
\noindent \textbf{University of Perpignan} \hfill  L3 Computer Science \\
 \hfill C. Legrand \hfill  Year \the\year \\
\smallskip
\begin{center}
{
  {\scshape \large Lab Work 2 - Security \\ Infection Techniques and Polymorphism}  \\
}
\end{center}
\hbox{\raisebox{0.4em}{\vrule depth 0pt height 0.9pt width \textwidth}}

\bigskip

\bigskip


\begin{exercice}[Logic Bomb with Temporal Trigger]

A \textit{logic bomb} is a malicious program that installs itself in a system and waits for a specific event called a \textbf{trigger} to execute its offensive function (the \textbf{payload}).

Consider the following C program:

\begin{verbatim}
#include <stdio.h>
#include <time.h>

int check_trigger() {
    time_t now = time(NULL);
    struct tm *t = localtime(&now);

    // Trigger condition
    if (t->tm_mon == 3 && t->tm_mday == 1) {
        return 1;
    }
    return 0;
}

void payload() {
    printf("=== PAYLOAD ACTIVATED ===\n");
    printf("All your files have been encrypted!\n");
}

void normal_operation() {
    printf("System running normally...\n");
}

int main() {
    if (check_trigger()) {
        payload();
    } else {
        normal_operation();
    }
    return 0;
}
\end{verbatim}

\begin{enumerate}
\item Identify the \textit{trigger} used in this program and explain precisely how it works. On which date will the payload be activated?

\item Modify the \texttt{check\_trigger()} function so that the trigger only activates if the program has been executed more than 100 times. \textit{Hint: use a file to store the execution counter.}

\item Propose two techniques to make the \textit{trigger} harder to detect during static code analysis.

\item From an antivirus perspective, what techniques could be used to detect this type of malicious program?
\end{enumerate}

\end{exercice}


\begin{exercice}[XOR Encryption and Polymorphism]

\textbf{Polymorphism} is a technique used by viruses to evade signature-based detection. By encrypting the viral code with different keys for each copy, the virus produces different binary signatures while maintaining the same behavior.

\textbf{XOR} (exclusive or) encryption is often used because it has an interesting property: applying the same operation twice with the same key returns the original value.

Consider the following C program:

\begin{verbatim}
#include <stdio.h>
#include <string.h>

void xor_crypt(char *data, int len, char key) {
    for (int i = 0; i < len; i++) {
        data[i] = data[i] ^ key;
    }
}

int main() {
    char payload[] = "FORMAT C: /Y";
    char key = 0x42;

    printf("Original : %s\n", payload);

    // Encryption
    xor_crypt(payload, strlen(payload), key);
    printf("Encrypted: ");
    for (int i = 0; i < strlen(payload); i++) {
        printf("%02X ", (unsigned char)payload[i]);
    }
    printf("\n");

    // Decryption
    xor_crypt(payload, strlen(payload), key);
    printf("Decrypted: %s\n", payload);

    return 0;
}
\end{verbatim}

\begin{enumerate}
\item Explain mathematically why the XOR operation allows both encryption and decryption using the same function. \textit{Hint: show that $A \oplus K \oplus K = A$.}

\item Mentally execute the program with key \texttt{0x42}. What is the first byte of the encrypted payload? (The ASCII code for 'F' is 0x46)

\item Modify the \texttt{xor\_crypt()} function to use a multi-byte key (e.g., \texttt{"SECRET"}) instead of a single byte.

\item From an antivirus perspective:
\begin{enumerate}
\item Why does this technique make signature-based detection ineffective?
\item What alternative technique could the antivirus use to detect this type of polymorphic code?
\end{enumerate}
\end{enumerate}

\end{exercice}


\begin{exercice}[Companion Virus and PATH Hijacking]

A \textbf{companion virus} is a type of malware that does not modify existing files but places itself alongside them to be executed in their place. A common technique involves exploiting the \texttt{PATH} environment variable to intercept the execution of system commands.

\textbf{Attack principle:} When a user types a command (for example \texttt{ls}), the system searches for the executable in the directories listed in \texttt{PATH}, in order. If an attacker places a fake \texttt{ls} in a higher-priority directory, this malicious program will be executed instead.

Consider the following C program, designed to replace the \texttt{ls} command:

\begin{verbatim}
#include <stdio.h>
#include <stdlib.h>
#include <unistd.h>

int main(int argc, char *argv[]) {
    // Phase 1: Malicious code (simulation)
    FILE *f = fopen("/tmp/log.txt", "a");
    if (f) {
        fprintf(f, "Intercepted command: ls\n");
        fprintf(f, "Directory: %s\n", getcwd(NULL, 0));
        fclose(f);
    }

    // Phase 2: Execute the real command
    execvp("/bin/ls", argv);

    return 0;
}
\end{verbatim}

\begin{enumerate}
\item Explain how this program works. Why doesn't the user notice that their system is compromised?

\item How could an attacker install this companion virus on a Linux system? Propose two different methods.

\item What happens if the attacker forgets the line \texttt{execvp("/bin/ls", argv)}? What consequences would this have on the stealthiness of the attack?

\item Propose two countermeasures that the user or system administrator could implement to protect against this type of attack.

\item \textbf{Bonus:} Modify the program to also record the arguments passed to the \texttt{ls} command.
\end{enumerate}

\end{exercice}


\begin{exercice}[Target File Search Routine]

Before being able to infect files, a virus must first locate them on the system. This \textbf{target search} phase is crucial: it must be efficient to maximize propagation while remaining stealthy to avoid detection.

Consider the following C program that recursively traverses a directory looking for files with a specific extension:

\begin{verbatim}
#include <stdio.h>
#include <stdlib.h>
#include <string.h>
#include <dirent.h>
#include <sys/stat.h>

void search_targets(const char *path, const char *extension) {
    DIR *dir = opendir(path);
    if (dir == NULL) return;

    struct dirent *entry;
    char fullpath[1024];

    while ((entry = readdir(dir)) != NULL) {
        // Ignore . and ..
        if (strcmp(entry->d_name, ".") == 0 ||
            strcmp(entry->d_name, "..") == 0)
            continue;

        snprintf(fullpath, sizeof(fullpath), "%s/%s", path, entry->d_name);

        struct stat st;
        if (stat(fullpath, &st) == -1) continue;

        if (S_ISDIR(st.st_mode)) {
            // Recurse into subdirectories
            search_targets(fullpath, extension);
        } else if (S_ISREG(st.st_mode)) {
            // Check the extension
            char *ext = strrchr(entry->d_name, '.');
            if (ext && strcmp(ext, extension) == 0) {
                printf("Target found: %s (%ld bytes)\n",
                       fullpath, st.st_size);
            }
        }
    }
    closedir(dir);
}

int main(int argc, char *argv[]) {
    if (argc != 3) {
        printf("Usage: %s <directory> <extension>\n", argv[0]);
        return 1;
    }

    printf("Searching for %s files in %s...\n", argv[2], argv[1]);
    search_targets(argv[1], argv[2]);

    return 0;
}
\end{verbatim}

\begin{enumerate}
\item Analyze the code and explain the role of the functions \texttt{opendir()}, \texttt{readdir()}, \texttt{stat()}, and \texttt{closedir()}.

\item Why is it important to ignore the "." and ".." entries when traversing a directory? What would happen otherwise?

\item Modify the \texttt{search\_targets()} function to search for executable files (having execution permission) instead of filtering by extension.
\textit{Hint: use \texttt{st.st\_mode \& S\_IXUSR} to test the execution permission.}

\item Recursive traversal can be problematic on large file systems. Propose a modification to limit the recursion depth to a maximum of $N$ levels.

\item From an antivirus perspective, what suspicious behaviors could this type of program trigger? How could a behavioral antivirus detect this activity?
\end{enumerate}

\end{exercice}


\begin{exercice}[Code Injection by Appending]

\textbf{Appending infection} is a classic technique where the virus adds its code to the end of a target file. To avoid infecting the same file multiple times (which could corrupt it or make it suspicious), the virus uses an \textbf{infection marker}.

Consider the following C program that simulates this mechanism on text files:

\begin{verbatim}
#include <stdio.h>
#include <string.h>
#include <stdlib.h>

#define MARKER "[[INFECTED]]"
#define PAYLOAD "\n--- MALICIOUS CODE SIMULATION ---\n"

int is_infected(const char *filename) {
    FILE *f = fopen(filename, "r");
    if (f == NULL) return -1;

    char buffer[1024];
    while (fgets(buffer, sizeof(buffer), f) != NULL) {
        if (strstr(buffer, MARKER) != NULL) {
            fclose(f);
            return 1;  // Already infected
        }
    }
    fclose(f);
    return 0;  // Not infected
}

int infect_file(const char *filename) {
    if (is_infected(filename) == 1) {
        printf("[SKIP] %s is already infected\n", filename);
        return 0;
    }

    FILE *f = fopen(filename, "a");
    if (f == NULL) {
        printf("[ERROR] Cannot open %s\n", filename);
        return -1;
    }

    fprintf(f, "%s", PAYLOAD);
    fprintf(f, "%s\n", MARKER);
    fclose(f);

    printf("[INFECTED] %s\n", filename);
    return 1;
}

int main(int argc, char *argv[]) {
    if (argc < 2) {
        printf("Usage: %s <file1> [file2] ...\n", argv[0]);
        return 1;
    }

    for (int i = 1; i < argc; i++) {
        infect_file(argv[i]);
    }
    return 0;
}
\end{verbatim}

\begin{enumerate}
\item Explain the role of the infection marker (\texttt{MARKER}). What would happen if the virus did not use a marker?

\item Why is the file opened in "a" (append) mode in the \texttt{infect\_file()} function? What would be the effect of using "w" mode instead?

\item The \texttt{is\_infected()} function scans the entire file line by line. This approach has performance limitations. Propose a more efficient method to check if a file is infected.
\textit{Hint: think about the position of the marker in the file.}

\item The marker \texttt{[[INFECTED]]} is easily detectable by an antivirus. Propose a modification to make the marker less visible while remaining functional for the virus.

\item From an antivirus perspective, what characteristics of this program could be used to detect it? Propose at least two detection techniques.
\end{enumerate}

\end{exercice}


\end{document}
